\subsection{后训练目标：会续写不等于会帮助用户}

slides 25--28 完成从 pretraining 到 post-training 的概念转折。语言模型目标只要求给训练分布中的文本高概率，不要求遵循当前用户意图、拒绝危险请求、使用指定格式或在困难任务上花更多 test-time compute。经典 alignment/post-training 以人类偏好和任务效用为目标；reasoning training 则强调有可验证答案的困难问题，并通过更长推理实现 test-time scaling。

\lecturefigure{slides-images/slide-025.png}{官方 slide 25：language modeling 与 assisting users 的目标错位}
\lecturefigure{slides-images/slide-026.png}{官方 slide 26：经典 post-training/alignment/instruction following}
\lecturefigure{slides-images/slide-027.png}{官方 slide 27：reasoning 训练的正确性目标与可验证任务}
\lecturefigure{slides-images/slide-028.png}{官方 slide 28：reasoning 与 test-time scaling}

\begin{knowledgebox}{术语消化：三个目标函数解决三类问题}
\begin{tabular}{p{2.5cm}p{4.2cm}p{4.5cm}}
\toprule
\textbf{阶段} & \textbf{优化信号} & \textbf{最直接解决的问题} \\
\midrule
Pretraining & next-token likelihood & 获得广泛知识、语言和模式先验 \\
Alignment/PT & 人类偏好或用户效用 & 指令遵循、风格、安全与交互习惯 \\
Reasoning RL & 可验证正确性奖励 & 在难题上探索多步策略与 test-time compute \\
\bottomrule
\end{tabular}
\end{knowledgebox}

\textbf{老师强调：}经典 post-training 可以理解为把“知道很多但不会与人交互”的模型变成 instruct model；reasoning RL 在真实训练顺序中常位于经典偏好对齐之前或与之交错。slide 3 把它们并列是为了比较，不是声明所有团队必须采用完全相同的固定顺序。对于 Agent，目标还要继续扩展到环境状态、工具调用和终局成功，因此后训练的定义本身会随任务边界变化。

